\section{Introduction}\label{sec:intro}

Board games have long served as a test of whether machines can learn and plan.
Programs now play checkers perfectly \parencite{schaeffer2007}, beat the best
humans at chess \parencite{campbell2002}, and have mastered Go by learning from
their own games \parencite{silver2016,silver2018}. International draughts, the
$10\times10$ form of draughts played mainly in Europe and Africa, has strong
programs too, the best freely available being Scan \parencite{letouzey_scan},
but it has received far less attention from research on learning.

This paper continues an earlier attempt of ours to teach a computer draughts with
a neural network. That attempt followed \textcite{chellapilla1999}, who evolved
networks for checkers through play alone, and it never reached a useful
strength. Here we take a more direct route. Instead of discovering good play from
scratch, our network learns from a teacher: it is trained to predict how Scan
judges millions of positions, and it is then used inside a fast alpha-beta
search, in the style of the efficiently updatable networks of computer shogi
\parencite{nasu2018}.

The engine is not only tested against other programs. Under the name valhalla it
plays rated games against people on the public server lidraughts.org, where it is
measured on the same rating scale as several thousand human players. Each game it
loses against a human, or draws against a much weaker one, becomes a lesson: the
positions where it went wrong are added to its training data, and a retrained
network replaces the current one if it plays at least as well.

We ask four questions. How strong is the engine against people, and how quickly
does its rating show this? At what stage of the game are its games won, drawn and
lost? How does it compare with its teacher? And does learning from individual
games make it stronger? In three days of blitz the engine played 127 games. It
won 71 and drew 5 of its 76 games against humans without a single loss, and after
only nine games its rating was higher than that of 99.9\% of the site's blitz
players. Against its teacher, however, it is still about 93 Elo weaker, and
against other programs nearly every game was drawn.

The rest of the paper is organised as follows. \Cref{sec:lit} reviews the history
of game-playing programs from the first ideas to learned evaluation functions and
international draughts, \cref{sec:method} describes the engine, its training and
how it plays and learns on lidraughts, \cref{sec:results} reports its games,
\cref{sec:disc} discusses what they show and where they fall short, and
\cref{sec:concl} concludes with directions for future work.
